\section{Future Work}
\label{sec:future}

Several directions remain open and we discuss each in turn.

\paragraph{Path-sensitive continuation analysis.}
The forward field-access analysis that drives schema self-healing 
is intentionally sound but path-insensitive. It collects every 
field projection on a bound variable that appears anywhere in the 
continuation, including projections inside branches that may never 
execute at runtime. This means the runtime can require the oracle 
to supply fields that are only accessed in unreachable code, 
forcing unnecessary retries. A path-sensitive analysis would track 
which projections are reachable under each possible branch 
condition and restrict the healing requirement accordingly. The 
main challenge is that branch conditions in \pact{} can depend on 
LLM-produced values, which are not known statically, so a 
path-sensitive treatment would need to account for the 
non-deterministic oracle in its reachability reasoning.

\paragraph{Richer type structure.}
The \pact{} type system is deliberately simple. Arrays are 
homogeneously typed, records check only declared fields with no 
row polymorphism or width subtyping, and there are no union or 
intersection types. Each of these restrictions was made to keep 
gradual validation tractable, but each also limits what workflows 
can express without routing values through \texttt{Schema} and 
deferring all shape checking to the post-materialization 
re-checker. Row polymorphism would allow records to carry 
additional fields beyond those declared, which is a natural fit 
for LLM-produced JSON that routinely includes extra metadata. 
Union types would allow a single \texttt{gen} binding to produce 
one of several shapes depending on the reasoning context, removing 
the need to route such cases through an untyped \texttt{Schema} 
binding. Refinement types would allow constraints on field values, 
not just field presence, which would make the re-checking 
discipline more expressive without changing its fundamental 
structure.

\paragraph{Effect tracking for host-language escapes.}
Inline \texttt{js} blocks execute arbitrary JavaScript under the 
host runtime's ambient authority. The current type system treats 
these blocks as oracular, accepting whatever value the JavaScript 
engine produces without modelling its side effects. This means 
that a \texttt{js} block can perform file I/O, network requests, 
or other host-side effects that are invisible to the \pact{} type 
checker and to the Cedar authorization model. An effect system 
that tracks which capabilities a \texttt{js} block exercises would 
allow the runtime to enforce that host-language escapes are 
consistent with the installed policy, extending the 
language-level control that \pact{} currently provides for 
\texttt{gen} and \texttt{agent} to the full program.

\begin{sloppypar}
\paragraph{Prompt-injection-aware information flow.}
The current threat model identifies prompt injection as outside
the scope of gradual validation. A \texttt{gen} prompt composed
from an \texttt{ask} result can be steered by an adversary who
controls the operator input, and the type system has no mechanism
to detect or prevent this. Information-flow tracking at the string
level would allow the runtime to label values by their source,
distinguishing programmer-authored text from operator-supplied
text from LLM-produced text, and to enforce policies over how
labeled values flow into prompt positions. The \texttt{ask} and
\texttt{gen} primitive separation already makes the source of each
value syntactically visible, which is a natural starting point
for such an analysis.
\end{sloppypar}

\paragraph{Parallel \texttt{gen} and modular workflow imports.}
The current operational semantics is fully sequential. All 
\texttt{gen} and \texttt{agent} calls execute one at a time under 
a single evolving Cedar policy. Many practical workflows could 
benefit from parallel generation, where independent \texttt{gen} 
calls that do not share data dependencies execute concurrently. 
Extending the semantics to support parallel \texttt{gen} would 
require a confluence argument showing that concurrent 
materialization and policy installation remain consistent, 
particularly given that \texttt{enforce} is append-only and 
policy monotonicity must be preserved across parallel steps. 
Separately, the current language has no module system. Workflows 
are declared within a single program and cannot be imported from 
external sources. Modular workflow imports would allow 
organizations to share typed workflow libraries, but would require 
an interface language for workflow types and a trust model for 
imported code that interacts cleanly with the existing Cedar 
authorization model.